Softwaretests sind essenziell für die Qualitätssicherung moderner Webanwendungen. Während die Testpyramide durch eine hohe Anzahl isolierter Unit-Tests eine gezielte Fehlerlokalisierung ermöglicht, setzt das Testing-Honeycomb-Modell auf realitätsnahe End-to-End-Tests, die den gesamten Anwendungsfluss über alle Schichten hinweg abdecken.

Diese Arbeit vergleicht beide Teststrategien anhand einer monolithischen Webanwendung zur Veranstaltungsorganisation. Dabei wurden für die Testpyramide sowohl eine H2-Datenbank als auch eine MySQL-Testcontainer-Umgebung genutzt, während das Testing-Honeycomb-Modell ausschließlich auf MySQL-Testcontainer setzte. Laufzeiten, Testabdeckung und Fehlerdiagnose wurden mit Spring Boot, Testcontainers und JaCoCo systematisch analysiert.

Die Ergebnisse zeigen, dass Tests in der Testpyramide durchschnittlich dreimal schneller ausgeführt werden, während das Testing-Honeycomb-Modell eine um 15 Prozent höhere Abdeckung der Systemprozesse erzielt. Die schichtweise Struktur der Testpyramide erlaubt eine präzisere Fehlerlokalisierung, während das Honeycomb-Modell den Debugging-Aufwand erhöht, da Fehler über alle Schichten hinweg analysiert werden müssen.

Ein weiterer Unterschied zeigt sich in der Testentwicklung mit Test-Driven Development (TDD). Beim Testing-Honeycomb-Modell können Tests bereits geschrieben werden, bevor Service- oder Repository-Klassen existieren, was einen schnelleren Einstieg in die Testentwicklung ermöglicht. Die Testpyramide erfordert hingegen eine schrittweise Implementierung der Schichten, bevor umfassende Tests durchgeführt werden können. Gleichzeitig ist das Setup im Testing-Honeycomb-Modell aufwendiger, da bereits zu Beginn Datenbankkonfigurationen und Testdaten berücksichtigt werden müssen.

Die Untersuchung zeigt, dass die Wahl der Teststrategie stark von den Anforderungen der Anwendung abhängt. Während die Testpyramide besonders für modulare, agil entwickelte Software geeignet ist, bietet das Testing-Honeycomb-Modell Vorteile für Systeme mit vielen interagierenden Komponenten. Ein hybrider Ansatz könnte beide Vorteile kombinieren, indem er Unit-Tests zur schnellen Validierung und gezielte End-to-End-Tests zur Systemabsicherung integriert.

Diese Arbeit liefert eine praxisnahe Evaluierung beider Teststrategien und zeigt auf, wie hybride Ansätze die Vorteile beider Modelle kombinieren könnten. Zukünftige Arbeiten sollten untersuchen, wie sich ein kombinierter Testansatz automatisiert implementieren lässt und welche Kombination für verschiedene Softwarearchitekturen optimal ist.

